\subsection{Normal Operator Consequences}

\begin{perspective} \label{per:normal-operators-coordinate-free-diagonal} \hlt{Normal Means Diagonal in Orthogonal Coordinates}\\
Complex spectral theorem states that normal operators are precisely the operators whose full structure is visible in an orthonormal eigenbasis. This is why normality is the natural coordinate-free version of being diagonal.
\end{perspective}

\begin{proposition} \label{prop:normal-polynomial-kernel-range-properties} \hlt{Basic Spectral Properties of Normal Operators}\\
Let $\mathcal{V}$ be a finite-dimensional complex inner-product space, let $T\in\LL(\mathcal{V})$ be normal, and let $p\in\mathbb{C}[z]$.
\begin{enumerate}[label=(\roman*)]
\setlength{\itemsep}{0pt}
\item $p(T)$ is normal and commutes with $T$.
\item $T^*$ is normal and $E_{\lambda}(T)=E_{\overline{\lambda}}(T^*)$ for every $\lambda\in\mathbb{C}$.
\item $\ker T=\ker T^*$ and $\operatorname{range}(T)=\operatorname{range}(T^*)$.
\end{enumerate}
\begin{proof}
Choose an orthonormal eigenbasis $\mathbf{e}_1,\ldots,\mathbf{e}_n$ with $T\mathbf{e}_j=\mu_j\mathbf{e}_j$.
\begin{enumerate}[label=(\roman*)]
\setlength{\itemsep}{0pt}
\item $p(T)\mathbf{e}_j=p(\mu_j)\mathbf{e}_j$, $p(T)$ is diagonal in this basis, hence normal. It commutes with $T$ by Proposition \ref{prop:operator-polynomial-products}.
\item Since $T^*\mathbf{e}_j=\overline{\mu_j}\mathbf{e}_j$, $T^*$ is diagonal in the same basis, hence normal. Also $T\mathbf{v}=\lambda\mathbf{v}$ iff $T^*\mathbf{v}=\overline{\lambda}\mathbf{v}$.
\item Part (ii) with $\lambda=0$ gives $\ker T=\ker T^*$. Both ranges are $\operatorname{span}\{\mathbf{e}_j:\mu_j\neq0\}$.
\end{enumerate}
\hlt{Structural Commentary}: Normality gives one orthonormal basis that diagonalises both $T$ and $T^*$.\\
\hlt{Key Insight}: Kernel and range are read from the nonzero diagonal entries in that basis.
\end{proof}
\end{proposition}

\begin{lemma} \label{lem:adjoints-of-finitely-many-normal-matrices-polynomial} \hlt{Adjoints of Normal Matrices as One Polynomial}\\
Let $A_1,\ldots,A_k$ be normal complex matrices. Then $\exists$ a polynomial $p\in\mathbb{C}[z]$ s.t. $p(A_j)=A_j^* \ \forall j=1,\ldots,k$.
\begin{proof}
Let $\Omega=\bigcup\limits_{j=1}^{k}\operatorname{spec}(A_j)$. Since $\Omega$ is finite, choose a Lagrange interpolation polynomial $p$ satisfying $p(\lambda)=\overline{\lambda}$ for every $\lambda\in\Omega$. For each $A_j$, Corollary \ref{cor:normal-spectral-resolution-functional-calculus} gives $p(A_j)=A_j^*$.\\
\hlt{Structural Commentary}: One polynomial interpolates conjugation on all spectra.\\
\hlt{Key Insight}: For normal matrices, $A^*$ is $p(A)$ with $p(\lambda)=\overline{\lambda}$.
\end{proof}
\end{lemma}

\begin{theorem} \label{thm:fuglede-putnam-normal-intertwining} \hlt{Fuglede--Putnam Theorem}\\
Let $A\in\operatorname{Mat}(m,m,\mathbb{C})$ and $B\in\operatorname{Mat}(n,n,\mathbb{C})$ be normal, and let $X\in\operatorname{Mat}(m,n,\mathbb{C})$.\\
Then $AX=XB$ $\Longleftrightarrow$ $A^*X=XB^*$.
\begin{proof}
\begin{enumerate}[label=\roman*.]
\setlength{\itemsep}{0pt}
\item[$(\Rightarrow)$] Assume $AX=XB$. Then $A^rX=XB^r$ for every $r\geq0$, so $p(A)X=Xp(B)$ for every polynomial $p$. By Lemma \ref{lem:adjoints-of-finitely-many-normal-matrices-polynomial}, choose $p$ with $p(A)=A^*$ and $p(B)=B^*$. Thus $A^*X=XB^*$.
\item[$(\Leftarrow)$] Assume $A^*X=XB^*$. Apply the same argument to the normal matrices $A^*$ and $B^*$, choosing $q$ with $q(A^*)=A$ and $q(B^*)=B$. Hence $AX=XB$.
\end{enumerate}
\hlt{Structural Commentary}: An intertwining relation passes through every polynomial in the two normal matrices.\\
\hlt{Key Insight}: Normality makes adjoints polynomial, so intertwining transfers to adjoints.
\end{proof}
\end{theorem}

\begin{method} \label{meth:finding-spectral-decomposition-normal-matrix} \hlt{Finding a Spectral Decomposition of a Normal Matrix}\\
For a normal matrix $A\in\operatorname{Mat}(n,n,\mathbb{C})$:
\begin{enumerate}[label=\roman*.]
\setlength{\itemsep}{0pt}
\item Find the distinct eigenvalues $\mu_1,\ldots,\mu_d$.
\item For each $j$, solve $(A-\mu_jI_n)\mathbf{x}=0$ to get $E_{\mu_j}(A)$.
\item Orthonormalise a basis of each eigenspace.
\item Concatenate these orthonormal bases as the columns of $U$; then $U$ is unitary and $A=U\Lambda U^*$, where $\Lambda$ repeats each $\mu_j$ according to $\dim E_{\mu_j}(A)$.
\end{enumerate}
\end{method}

\begin{theorem} \label{thm:normal-spectral-classification-by-eigenvalues} \hlt{Spectral Classification of Normal Operators}\\
Let $\mathcal{V}$ be a finite-dimensional complex inner-product space and let $T\in\LL(\mathcal{V})$ be normal.
\begin{enumerate}[label=(\roman*)]
\setlength{\itemsep}{0pt}
\item $T$ is unitary if and only if $\abs{\lambda}=1$ for every $\lambda\in\operatorname{spec}(T)$.
\item $T$ is self-adjoint if and only if $\operatorname{spec}(T)\subseteq\mathbb{R}$.
\item $T$ is an orthogonal projection if and only if $\operatorname{spec}(T)\subseteq\{0,1\}$.
\item $T$ is skew-adjoint, meaning $T^*=-T$, if and only if every $\lambda\in\operatorname{spec}(T)$ is purely imaginary.
\end{enumerate}
\begin{proof}
Write the spectral resolution $T=\sum\limits_j\mu_jP_j$, where the $\mu_j$ are the distinct eigenvalues and $P_j\neq0$.\\
Then $I_{\mathcal{V}}=\sum_jP_j$, $T^*=\sum\limits_j\overline{\mu_j}P_j$, $T^*T=\sum\limits_j\abs{\mu_j}^2P_j$, and $T^2=\sum\limits_j\mu_j^2P_j$.\\
Since ranges of $P_j$ are nonzero and independent, an identity $\sum_j\alpha_jP_j=0$ holds only when every $\alpha_j=0$.
\begin{enumerate}[label=(\roman*)]
\setlength{\itemsep}{0pt}
\item $T$ is unitary iff $T^*T=I_{\mathcal{V}}=\sum_jP_j$, iff $\abs{\mu_j}=1$ for every $j$.
\item $T$ is self-adjoint iff $T=T^*$, iff $\mu_j=\overline{\mu_j}$ for every $j$.
\item $T$ is an orthogonal projection iff $T=T^*$ and $T^2=T$, iff $\mu_j=\overline{\mu_j}$ and $\mu_j^2=\mu_j \ \forall j$, iff $\mu_j\in\{0,1\} \ \forall j$.
\item $T^*=-T$ iff $\overline{\mu_j}=-\mu_j$ for every $j$, iff every $\mu_j$ is purely imaginary.
\end{enumerate}
\hlt{Structural Commentary}: For normal operators, structural classes are scalar conditions on spectral blocks.\\
\hlt{Key Insight}: For normal $T$, these properties reduce to scalar tests.
\end{proof}
\end{theorem}

\begin{proposition} \label{prop:cartesian-decomposition-normal-commuting-parts} \hlt{Cartesian Decomposition and Normality}\\
Let $\mathcal{V}$ be a finite-dimensional complex inner-product space.\\
For every $T\in\LL(\mathcal{V})$, there are unique self-adjoint operators $H,K\in\LL(\mathcal{V})$ such that $T=H+iK$.\\
They are $H=\frac{1}{2}(T+T^*)$ and $K=\frac{1}{2i}(T-T^*)$. Moreover, $T$ is normal if and only if $HK=KH$.
\begin{proof}
The displayed formulas give self-adjoint $H,K$ and $T=H+iK$. If $T=X+iY$ with $X,Y$ self-adjoint, then $T^*=X-iY$, so $X=\frac{1}{2}(T+T^*)$ and $Y=\frac{1}{2i}(T-T^*)$; uniqueness follows.\\
For $T=H+iK$, $T^*T=(H-iK)(H+iK)=H^2+K^2+i(HK-KH)$,\\
whereas $TT^*=(H+iK)(H-iK)=H^2+K^2-i(HK-KH)$. Thus $T^*T=TT^*$ if and only if $HK=KH$.\\
\hlt{Structural Commentary}: Every complex operator splits into Hermitian and skew-Hermitian parts.\\
\hlt{Key Insight}: Normality says the two Hermitian coordinates of $T$ commute.
\end{proof}
\end{proposition}

\begin{figure}[H]
  \centering
  \begin{tikzpicture}[
  >=Stealth,
  every node/.style={font=\small},
  box/.style={draw=gray!55, rounded corners=3pt, fill=gray!6, align=center, inner xsep=7pt, inner ysep=4.5pt, minimum height=0.9cm},
  result/.style={box, fill=green!7, minimum width=4.05cm},
  arrow/.style={->, line width=0.85pt, blue!70!black}
]
  \node[box, minimum width=1.4cm] (T) at (0,0) {$T$};
  \node[box, minimum width=3.75cm] (parts) at (2.95,0)
    {$H=\frac12(T+T^*)$\\[-1pt]
     $K=\frac{1}{2i}(T-T^*)$\\[-1pt]
     $T=H+iK$};
  \node[result] (normal) at (7.35,0)
    {$T\ \text{normal}\quad\Longleftrightarrow\quad HK=KH$};

  \draw[arrow] (T) -- (parts);
  \draw[arrow] (parts) -- (normal);
\end{tikzpicture}

  \caption{Cartesian decomposition: $T$ is normal exactly when $HK=KH$.}
  \label{fig:cartesian-normal-commutation}
\end{figure}

\begin{proposition} \label{prop:block-upper-triangular-normal-collapse} \hlt{Normal Block Upper-Triangular Matrices Collapse}\\
Let $\mathbb{F}\in\{\mathbb{R},\mathbb{C}\}$. Let $A=
\begin{bmatrix}
B&X\\
0&C
\end{bmatrix}
\in\operatorname{Mat}(n,n,\mathbb{F})$, where $B$ and $C$ are square.\\
Then $A$ is normal if and only if $X=0$ and both $B$ and $C$ are normal.
\begin{proof}
\begin{enumerate}[label=\roman*.]
\setlength{\itemsep}{0pt}
\item[$(\Rightarrow)$] Suppose $A$ is normal. Comparing the $(1,1)$ blocks of $A^*A$ and $AA^*$ gives $B^*B=BB^*+XX^*$.\\
Taking traces and using cyclicity, $\operatorname{tr}(XX^*)=\operatorname{tr}(B^*B)-\operatorname{tr}(BB^*)=0$. Since $\operatorname{tr}(XX^*)=\norm{X}_F^2$, we get $X=0$. Thus $A=B\oplus C$, and $A^*A=AA^*$ reduces blockwise to $B^*B=BB^*$ and $C^*C=CC^*$.
\item[$(\Leftarrow)$] If $X=0$ and $B,C$ are normal, then $A=B\oplus C$, so blockwise multiplication gives $A^*A=AA^*$.
\end{enumerate}
\hlt{Structural Commentary}: A normal operator cannot have one-way coupling between orthogonal blocks.\\
\hlt{Key Insight}: The off-diagonal block contributes the nonnegative trace $\norm{X}_F^2$, so normality forces it to vanish.
\end{proof}
\end{proposition}

\begin{figure}[H]
  \centering
  \begin{tikzpicture}[
  >=Stealth,
  every node/.style={font=\small},
  block/.style={draw=gray!55, rounded corners=3pt, inner xsep=6pt, inner ysep=5pt, align=center},
  arrow/.style={->, line width=0.85pt, blue!70!black}
]
  \node[block] (tri) at (0,0)
    {$A=\begin{bmatrix}B&X\\0&C\end{bmatrix}$\\normal};
  \node[block] (schur) at (3.55,0)
    {$T_B=U_B^*BU_B,\quad T_C=U_C^*CU_C$\\
     $Q^*AQ=\begin{bmatrix}T_B&U_B^*XU_C\\0&T_C\end{bmatrix}$\\
     upper triangular and normal};
  \node[block] (diag) at (7.62,0)
    {$Q^*AQ$ diagonal\\$X=0$, $B,C$ normal};

  \draw[arrow] (tri) -- (schur);
  \draw[arrow] (schur) -- (diag);
\end{tikzpicture}

  \caption{Normality removes the upper-right block in a block upper-triangular matrix.}
  \label{fig:block-upper-triangular-normal-collapse}
\end{figure}

\begin{theorem} \label{thm:normal-invariant-subspaces-reduce} \hlt{Invariant Subspaces of Normal Matrices Reduce}\\
Let $\mathbb{F}\in\{\mathbb{R},\mathbb{C}\}$, let $A\in\operatorname{Mat}(n,n,\mathbb{F})$ be normal, and let $\mathcal{U}\subseteq\mathbb{F}^n$ be invariant under $A$. Then $\mathcal{U}$ is invariant under $A^*$. If $P$ is the orthogonal projection onto $\mathcal{U}$, then $AP=PA$.
\begin{proof}
The cases $\mathcal{U}=\{0\}$ and $\mathcal{U}=\mathbb{F}^n$ are immediate. Otherwise choose a unitary $Q=[Q_1\ Q_2]$ whose columns form orthonormal bases of $\mathcal{U}$ and $\mathcal{U}^{\perp}$. Invariance of $\mathcal{U}$ says $AQ_1$ has columns in $\mathcal{U}=\operatorname{col}(Q_1)$, hence $Q_2^*AQ_1=0$.\\
Therefore $Q^*AQ=
\begin{bmatrix}
B&X\\
0&C
\end{bmatrix}$. This matrix is normal because it is unitarily similar to $A$.\\
Proposition \ref{prop:block-upper-triangular-normal-collapse} gives $X=0$, so $Q^*AQ=B\oplus C$. Hence $Q^*A^*Q=B^*\oplus C^*$, which shows $\mathcal{U}$ is $A^*$-invariant.\\
Also $Q^*PQ=I_{\dim(\mathcal{U})}\oplus0$, so $Q^*AQ$ commutes with $Q^*PQ$; conjugating back gives $AP=PA$.\\
\hlt{Structural Commentary}: In adapted coordinates, normal invariance becomes block diagonality.\\
\hlt{Key Insight}: Normal invariant subspaces are automatically reducing subspaces.
\end{proof}
\end{theorem}

\begin{corollary} \label{cor:normal-operator-invariant-subspaces-reduce} \hlt{Invariant Subspaces of Normal Operators Reduce}\\
Let $\mathcal{V}$ be a finite-dimensional inner-product space, let $T\in\LL(\mathcal{V})$ be normal, let $\mathcal{U}$ be invariant under $T$. Then:
\begin{enumerate}[label=(\roman*)]
\setlength{\itemsep}{0pt}
\item $\mathcal{U}^{\perp}$ is invariant under $T$;
\item $\mathcal{U}$ is invariant under $T^*$;
\item $(T|_{\mathcal{U}})^*=(T^*)|_{\mathcal{U}}$;
\item $T|_{\mathcal{U}}$ and $T|_{\mathcal{U}^{\perp}}$ are normal.
\end{enumerate}
\begin{proof}
Choose an orthonormal basis adapted to $\mathcal{U}\oplus\mathcal{U}^{\perp}$.\\
In these coordinates, Theorem \ref{thm:normal-invariant-subspaces-reduce} makes the matrix of $T$ block diagonal, say $A=B\oplus C$.\\
The block diagonal form shows $\mathcal{U}^{\perp}$ is invariant under $T$ and $\mathcal{U}$ is invariant under $T^*$. It also gives $(T|_{\mathcal{U}})^*=(T^*)|_{\mathcal{U}}$ because the adjoint block on $\mathcal{U}$ is $B^*$. Finally, normality of $A$ reduces blockwise to normality of $B$ and $C$.\\
\hlt{Structural Commentary}: Normal operator turns every invariant subspace into an orthogonal reducing summand.\\
\hlt{Key Insight}: Normality kills the one-way coupling between $\mathcal{U}$ and $\mathcal{U}^{\perp}$.
\end{proof}
\end{corollary}

\begin{perspective} \label{per:real-normal-operators-planes} \hlt{Real Normal Operators Use Planes}\\
Over $\mathbb{R}$, normal operators need not have enough eigenvectors. The replacement for a missing conjugate eigenpair is a two-dimensional invariant plane carrying a rotation-dilation block.
\end{perspective}

\begin{proposition} \label{prop:two-dimensional-real-normal-block} \hlt{Two-Dimensional Real Normal Blocks}\\
Let $\dim(\mathcal{V})=2$ over $\mathbb{R}$ and let $T\in\LL(\mathcal{V})$.\\
Then $T$ is normal but not self-adjoint if and only if, in every orthonormal basis, $[T]=
\begin{bmatrix}
a&-b\\
b&a
\end{bmatrix}$, $b\neq0$.\\
Moreover, some orthonormal basis makes $b>0$.
\begin{proof}
\begin{enumerate}[label=\roman*.]
\setlength{\itemsep}{0pt}
\item[$(\Rightarrow)$] In an orthonormal basis, write $A=[T]=
\begin{bmatrix}
a&c\\
b&d
\end{bmatrix}$.\\
Normality gives $A^tA=AA^t$. Comparing the $(1,1)$ and $(2,2)$ entries gives $b^2=c^2$. If $c=b$, then $A=A^t$, impossible because $T$ is not self-adjoint. Thus $c=-b$ and $b\neq0$. Comparing off-diagonal entries gives $-ab+bd=ab-bd$, so $d=a$. Since the original orthonormal basis was arbitrary, every orthonormal basis has the displayed form.
\item[$(\Leftarrow)$] A matrix $\begin{bmatrix}a&-b\\ b&a\end{bmatrix}$ with $b\neq0$ satisfies $A^tA=AA^t$ and is not symmetric. Hence $T$ is normal but not self-adjoint. If $b<0$, replace the second basis vector by its negative to make $b>0$.
\end{enumerate}
\hlt{Structural Commentary}: In real dimension $2$, non-self-adjoint normality is exactly a scalar plus a rotation block.\\
\hlt{Key Insight}: Normality forces $c=-b$ and $d=a$.
\end{proof}
\end{proposition}

\begin{theorem} \label{thm:real-normal-operator-block-diagonal-form} \hlt{Real Normal Operator Block Form}\\
Let $\mathcal{V}$ be finite-dimensional over $\mathbb{R}$ and let $T\in\LL(\mathcal{V})$. Then $T$ is normal if and only if some orthonormal basis makes $[T]$ block diagonal, with each diagonal block either $1$-by-$1$ or of the form $\begin{bmatrix}
a&-b\\
b&a
\end{bmatrix}$, $b > 0$.
\begin{proof}
\begin{enumerate}[label=\roman*.]
\setlength{\itemsep}{0pt}
\item[$(\Rightarrow)$] Induct on $n=\dim(\mathcal{V})$. The cases $n\leq1$ are immediate.\\
For $n\geq2$, Theorem \ref{thm:real-operator-invariant-line-or-plane} gives an invariant subspace $\mathcal{U}$ with $\dim(\mathcal{U})\in\{1,2\}$.\\
By Corollary \ref{cor:normal-operator-invariant-subspaces-reduce}, $\mathcal{U}^{\perp}$ is invariant under $T$ and both restrictions are normal.\\
If $\dim(\mathcal{U})=1$, take a unit vector in $\mathcal{U}$. If $\dim(\mathcal{U})=2$, use the Real Spectral Theorem \ref{thm:real-spectral-theorem} when $T|_{\mathcal{U}}$ is self-adjoint, giving scalar $1$-by-$1$ blocks on $\mathcal{U}$, and use Proposition \ref{prop:two-dimensional-real-normal-block} otherwise.\\
By induction, choose the required orthonormal block basis on $\mathcal{U}^{\perp}$. Concatenating it with the basis chosen on $\mathcal{U}$ gives an orthonormal basis of $\mathcal{V}$ because $\mathcal{V}=\mathcal{U}\oplus\mathcal{U}^{\perp}$.
\item[$(\Leftarrow)$] Each allowed block is normal. If $[T]$ is the block diagonal sum of these blocks, then $[T]^*[T]$ and $[T][T]^*$ are the corresponding block diagonal sums of $B^*B$ and $BB^*$, hence are equal block by block.
\end{enumerate}
\hlt{Structural Commentary}: Normality lets invariant lines or planes split orthogonally.\\
\hlt{Key Insight}: Real normal operators split into scalar lines and rotation-dilation planes.
\end{proof}
\end{theorem}

\begin{figure}[H]
  \centering
  \begin{tikzpicture}[
  >=Stealth,
  every node/.style={font=\small},
  box/.style={draw=gray!55, rounded corners=3pt, inner xsep=8pt, inner ysep=6pt, align=center},
  arrow/.style={->, line width=0.85pt, blue!70!black}
]
  \node[box] (normal) at (0,0) {normal $T$};
  \node[box] (sum) at (3.6,0) {$\mathcal{V}=\bigoplus L_j\oplus\bigoplus P_k$};
  \node[box] (blocks) at (7.8,0) {$[T]=\operatorname{diag}(\lambda_j,R_k)$\\[-1pt]
  $R_k=\begin{bmatrix}a_k&-b_k\\ b_k&a_k\end{bmatrix}$};

  \draw[arrow] (normal.east) -- (sum.west);
  \draw[arrow] (sum.east) -- (blocks.west);
\end{tikzpicture}

  \caption{Orthogonal line/plane split into scalar and rotation blocks.}
  \label{fig:real-normal-block-structure}
\end{figure}

\begin{theorem} \label{thm:simultaneous-unitary-diagonalization-normal-family} \hlt{Simultaneous Unitary Diagonalisation of Commuting Normal Operators}\\
Let $\emptyset\neq\mathcal{F}\subseteq\LL(\mathcal{V})$ be a normal family on a finite-dimensional complex inner-product space. Then $\mathcal{F}$ is pairwise commuting if and only if $\mathcal{V}$ has an orthonormal basis of common eigenvectors for $\mathcal{F}$.
\begin{proof}
\begin{enumerate}[label=\roman*.]
\setlength{\itemsep}{0pt}
\item[$(\Leftarrow)$] In a common orthonormal eigenbasis, every $A\in\mathcal{F}$ is diagonal. Hence all operators in $\mathcal{F}$ commute.
\item[$(\Rightarrow)$] Induct on $\dim(\mathcal{V})$. The zero-dimensional case is vacuous. If every $A\in\mathcal{F}$ is scalar, any orthonormal basis works. Otherwise choose a nonscalar $A_0\in\mathcal{F}$ and write its spectral resolution as $A_0=\sum_j\mu_jP_j$.\\
Since $A_0$ is nonscalar, at least two nonzero ranges $\operatorname{range}(P_j)$ occur, and each has dimension strictly smaller than $\dim(\mathcal{V})$. By Corollary \ref{cor:normal-spectral-resolution-functional-calculus}, each $P_j$ is a polynomial in $A_0$.\\
If $B\in\mathcal{F}$, then $BA_0=A_0B$, so $B$ commutes with every polynomial in $A_0$, in particular with $P_j$. Hence $B(\operatorname{range}(P_j))\subseteq\operatorname{range}(P_j)$ because $P_jB\mathbf{x}=BP_j\mathbf{x}=B\mathbf{x}$ whenever $P_j\mathbf{x}=\mathbf{x}$.\\
The same commutation preserves $\ker(P_j)=\operatorname{range}(P_j)^{\perp}$.\\
In the decomposition $\mathcal{V}=\operatorname{range}(P_j)\oplus\ker(P_j)$, the matrix of $B$ is therefore block diagonal; since $B$ is normal, its diagonal restrictions are normal block by block.\\
Thus, on each nonzero $\operatorname{range}(P_j)$, the restricted family is again pairwise commuting and normal.\\
By induction it has a common orthonormal eigenbasis. Concatenating these bases over all $j$ gives a common orthonormal eigenbasis of $\mathcal{V}$.
\end{enumerate}
\hlt{Structural Commentary}: Spectral projections reduce the whole commuting family.\\
\hlt{Key Insight}: Diagonalise one nonscalar member, then induct on its eigenspaces.
\end{proof}
\end{theorem}

\begin{figure}[H]
  \centering
  \begin{tikzpicture}[
  >=Stealth,
  every node/.style={font=\small},
  block/.style={draw=gray!55, rounded corners=3pt, inner xsep=8pt, inner ysep=4.5pt, align=center},
  arrow/.style={->, line width=0.8pt, blue!70!black, shorten >=2pt, shorten <=2pt}
]
  \node[block] (family) at (0,0) {$\mathcal{F}$};
  \node[block] (split) at (2.46,0) {$A_0=\sum_j\mu_jP_j$};
  \node[block] (pieces) at (5.32,0) {$\bigoplus_j\operatorname{range}(P_j)$};
  \node[block] (basis) at (8.49,0) {common eigenbasis};

  \draw[arrow] (family.east) -- (split.west);
  \draw[arrow] (split.east) -- (pieces.west);
  \draw[arrow] (pieces.east) -- (basis.west);
\end{tikzpicture}

  \caption{For $A_0=\sum_j\mu_jP_j$, commutativity gives $BP_j=P_jB$; the ranges reduce $\mathcal{F}$, and induction gives a common orthonormal eigenbasis.}
  \label{fig:simultaneous-normal-diagonalization}
\end{figure}

\begin{corollary} \label{cor:commuting-normal-sums-products-normal} \hlt{Sums and Products of Commuting Normal Operators}\\
Let $S,T\in\LL(\mathcal{V})$ be commuting normal operators on a finite-dimensional complex inner-product space.\\
Then $S+T$ and $ST$ are normal.
\begin{proof}
By Theorem \ref{thm:simultaneous-unitary-diagonalization-normal-family}, $S$ and $T$ are diagonal in a common orthonormal basis.\\
Then $S+T$ and $ST$ are also diagonal in that basis, hence normal.\\
\hlt{Structural Commentary}: Commuting normal operators share orthonormal coordinates.\\
\hlt{Key Insight}: In common diagonal form, sums and products remain diagonal.
\end{proof}
\end{corollary}

\begin{theorem} \label{thm:complex-isometry-description} \hlt{Description of Complex Isometries}\\
Let $\mathcal{V}$ be a finite-dimensional complex inner-product space and let $S\in\LL(\mathcal{V})$. Then $S$ is an isometry if and only if $\mathcal{V}$ has an orthonormal basis consisting of eigenvectors of $S$ whose eigenvalues all have absolute value $1$.
\begin{proof}
If $S\mathbf{e}_j=\lambda_j\mathbf{e}_j$ in an orthonormal basis with $\abs{\lambda_j}=1$, then for $\mathbf{v}=\sum_j c_j\mathbf{e}_j$,\\
$\norm{S\mathbf{v}}^2=\sum\limits_j\abs{\lambda_jc_j}^2=\sum\limits_j\abs{c_j}^2=\norm{\mathbf{v}}^2$, so $S$ is an isometry. Conversely, if $S$ is an isometry, then $S^*S=I_{\mathcal{V}}$ by Theorem \ref{thm:finite-dimensional-isometry-iff-adjoint-product-identity}. Finite-dimensional invertibility gives $SS^*=I_{\mathcal{V}}$, so $S$ is normal. By Complex Spectral Theorem, choose orthonormal eigenbasis $S\mathbf{e}_j=\lambda_j\mathbf{e}_j$. Since $\mathbf{e}_j$ is unit and $S$ preserves norms, $\abs{\lambda_j}=\norm{S\mathbf{e}_j}=1$.\\
\hlt{Structural Commentary}: Complex isometries are normal with unit-circle spectrum.\\
\hlt{Key Insight}: In an eigenbasis, isometry means $\abs{\lambda_j}=1$.
\end{proof}
\end{theorem}

\begin{example} \label{ex:plane-rotation-isometry} \hlt{Plane Rotations}\\
For $\theta\in\mathbb{R}$, the real matrix $R_{\theta}=
\begin{bmatrix}
\cos\theta&-\sin\theta\\
\sin\theta&\cos\theta
\end{bmatrix}$ is an isometry of $\mathbb{R}^2$. If $\theta$ is not an integer multiple of $\pi$, then $R_{\theta}$ has no real eigenvalue, which is why real isometries require two-dimensional rotation blocks.
\end{example}

\begin{theorem} \label{thm:real-isometry-block-diagonal-form} \hlt{Real Isometry Block Form}\\
Let $\mathcal{V}$ be finite-dimensional over $\mathbb{R}$ and let $S\in\LL(\mathcal{V})$. Then $S$ is an isometry if and only if some orthonormal basis makes $[S]$ block diagonal, with each diagonal block one of $[1]$, $[-1]$, $\begin{bmatrix}
\cos\theta&-\sin\theta\\
\sin\theta&\cos\theta
\end{bmatrix}$ with $0<\theta<\pi$.
\begin{proof}
\begin{enumerate}[label=\roman*.]
\setlength{\itemsep}{0pt}
\item[$(\Rightarrow)$] Since $S$ is an isometry, $S^*S=I_{\mathcal{V}}$ by Theorem \ref{thm:finite-dimensional-isometry-iff-adjoint-product-identity}. Finite-dimensional invertibility gives $SS^*=I_{\mathcal{V}}$, so $S$ is normal. Apply Theorem \ref{thm:real-normal-operator-block-diagonal-form}. A $1$-by-$1$ block $[\lambda]$ has $\lambda^2=1$, so $\lambda=\pm1$.\\
A $2$-by-$2$ block $\begin{bmatrix}a&-b\\b&a\end{bmatrix}$ has $a^2+b^2=1$ and $b>0$, so $a=\cos\theta$ and $b=\sin\theta$ for a unique $\theta\in(0,\pi)$.
\item[$(\Leftarrow)$] Each displayed block $B$ satisfies $B^tB=I$ on its coordinate subspace. Therefore the block diagonal matrix satisfies $[S]^t[S]=I$, so $S$ preserves norms and is an isometry.
\end{enumerate}
\hlt{Structural Commentary}: Isometries are the unit-length special case of real normal blocks.\\
\hlt{Key Insight}: Scalar blocks reduce to $\pm1$, and planar normal blocks reduce to rotations.
\end{proof}
\end{theorem}

\begin{figure}[H]
  \centering
  \begin{tikzpicture}[
  >=Stealth,
  every node/.style={font=\small},
  box/.style={draw=gray!55, rounded corners=3pt, inner xsep=8pt, inner ysep=6pt, align=center},
  arrow/.style={->, line width=0.85pt, blue!70!black}
]
  \node[box] (iso) at (0,0) {isometry $S$};
  \node[box] (blocks) at (3.3,0) {orthogonal sum\\$[1]$, $[-1]$, $R_{\theta}$};
  \node[box] (rot) at (7.2,0) {$R_{\theta}=\begin{bmatrix}\cos\theta&-\sin\theta\\ \sin\theta&\cos\theta\end{bmatrix}$};

  \draw[arrow] (iso) -- (blocks);
  \draw[arrow] (blocks) -- (rot);
\end{tikzpicture}

  \caption{Normal form: fixed/flipped lines and rotation planes.}
  \label{fig:real-isometry-blocks}
\end{figure}
